%! AP Statistics — Unit 1, Lesson 1.5 — Homework
\documentclass[10pt]{article}
\usepackage{apstats-article}
\usepackage{apstats-boxes}

%=============================================================================
\begin{document}

\pageheader{Unit 1, Lesson 1.5}{Homework: Representing a Quantitative Variable with Graphs}
\namedateperiod

\vspace{0.1in}

\begin{tcolorbox}[colback=sky, colframe=navy, boxrule=0.8pt, arc=2mm,
    left=3mm, right=3mm, top=1.5mm, bottom=1.5mm]
\small For all problems: show calculations, state context in written responses,
and use complete sentences for AP-style questions.
\end{tcolorbox}

\vspace{0.12in}

% ════════════════════════════════════════════════════════════════════════════
% PART A
% ════════════════════════════════════════════════════════════════════════════
\begin{blurbbox}[Part A \quad Interpret a Dotplot]{navylight}
\small

The dotplot below shows the time (minutes) it took each of 20 students to
complete an in-class problem set.

\vspace{0.12in}

\begin{center}
\begin{tikzpicture}[scale=1.05]
  \draw[thick] (4.2,0) -- (11.8,0);
  \foreach \x/\lbl in {5/5, 6/6, 7/7, 8/8, 9/9, 10/10, 11/11} {
    \draw (\x,-0.12) -- (\x,0.12);
    \node[below, font=\small] at (\x,-0.12) {\lbl};
  }
  \node[below, font=\small\itshape] at (8,-0.52) {Time (minutes)};
  % freq: 5→1, 6→3, 7→5, 8→5, 9→3, 10→2, 11→1
  \filldraw[navy] (5, 0.38) circle (0.14);
  \foreach \k in {1,...,3} \filldraw[navy] (6, \k*0.38) circle (0.14);
  \foreach \k in {1,...,5} \filldraw[navy] (7, \k*0.38) circle (0.14);
  \foreach \k in {1,...,5} \filldraw[navy] (8, \k*0.38) circle (0.14);
  \foreach \k in {1,...,3} \filldraw[navy] (9, \k*0.38) circle (0.14);
  \foreach \k in {1,...,2} \filldraw[navy] (10, \k*0.38) circle (0.14);
  \filldraw[navy] (11, 0.38) circle (0.14);
\end{tikzpicture}
\end{center}

\vspace{0.12in}

\begin{enumerate}[leftmargin=1.5em, itemsep=8pt]

  \item What is the range of completion times? \blank{3cm}
        How many students finished in 8 minutes or less?
        \blank{2.5cm}

  \item Describe the shape of this distribution. Is it symmetric, skewed left,
        or skewed right? Explain your reasoning.\\[6pt]
        \writeline\\[1pt]\writeline

  \item What percentage of students took longer than 9 minutes to finish?
        Show your calculation.\\[6pt]
        \writeline

  \item Write 2--3 complete sentences describing the distribution of completion
        times for these 20 students. Include shape, a typical value, and spread.\\[6pt]
        \writeline\\[1pt]\writeline\\[1pt]\writeline

\end{enumerate}

\end{blurbbox}

\vspace{0.12in}

% ════════════════════════════════════════════════════════════════════════════
% PART B
% ════════════════════════════════════════════════════════════════════════════
\begin{blurbbox}[Part B \quad Complete a Stemplot]{greenacc}
\small

Daily high temperatures ($^\circ$F) were recorded for 15 days in October:
\[52,\;55,\;58,\;61,\;63,\;65,\;67,\;68,\;70,\;72,\;73,\;74,\;76,\;78,\;82\]

The stems have been provided below. Complete the stemplot by filling in the leaves
for each stem. Then answer the questions.

\vspace{0.1in}

\begin{center}
\renewcommand{\arraystretch}{1.6}
\begin{tabular}{r | l}
\textbf{Stem} & \textbf{Leaf} \\
\hline
5 & \blank{5cm} \\
6 & \blank{5cm} \\
7 & \blank{5cm} \\
8 & \blank{5cm} \\
\end{tabular}
\hspace{1.2cm} Key: \textbf{7 $|$ 0} = 70$^\circ$F
\end{center}

\vspace{0.1in}

\begin{enumerate}[leftmargin=1.5em, itemsep=8pt]

  \item What is the median temperature? Show how you used the stemplot to find it.\\[6pt]
        \writeline\\[1pt]\writeline

  \item Describe the shape of the distribution. Does it appear skewed or symmetric?
        What does that tell you about temperatures in October for this location?\\[6pt]
        \writeline\\[1pt]\writeline

  \item What is the interquartile range (IQR)? (Hint: find Q1 and Q3 first.)
        Show your work.\\[6pt]
        \writeline\\[1pt]\writeline

\end{enumerate}

\end{blurbbox}

\vspace{0.12in}

% ════════════════════════════════════════════════════════════════════════════
% PART C
% ════════════════════════════════════════════════════════════════════════════
\begin{blurbbox}[Part C \quad AP-Style: Interpret a Histogram]{redacc}
\small

The frequency table below shows the distribution of the number of minutes 40 high
school students spent on homework on a given night.

\vspace{0.1in}

\renewcommand{\arraystretch}{1.6}
\begin{tabularx}{0.75\linewidth}{@{} C{0.26\linewidth} C{0.22\linewidth} C{0.26\linewidth} @{}}
\rowcolor{sky}
\textbf{Time (min)} & \textbf{Frequency} & \textbf{Rel.\ Frequency} \\
\midrule
$[0, 10)$   & 2  & \blank{2.5cm} \\
$[10, 20)$  & 8  & \blank{2.5cm} \\
$[20, 30)$  & 14 & \blank{2.5cm} \\
$[30, 40)$  & 11 & \blank{2.5cm} \\
$[40, 50)$  & 5  & \blank{2.5cm} \\
\midrule
\rowcolor{sky}\textbf{Total} & 40 & \blank{2.5cm} \\
\end{tabularx}

\vspace{0.12in}

\begin{enumerate}[leftmargin=1.5em, itemsep=10pt]

  \item Complete the relative frequency column. Verify that the relative
        frequencies sum to 1.

  \item Which bin contains the most students? What percentage of students
        fall in this bin?\\[6pt]
        \writeline

  \item What percentage of students spent at least 30 minutes on homework?
        Show your work.\\[6pt]
        \writeline

  \item Write a complete AP-style description of the distribution of homework
        time for these students. Address shape, center, and spread.\\[6pt]
        \writeline\\[1pt]\writeline\\[1pt]\writeline

  \item A classmate says: ``A stemplot would have been a better choice for this
        data.'' Do you agree? Justify your answer.\\[6pt]
        \writeline\\[1pt]\writeline

\end{enumerate}

\end{blurbbox}

\vspace{0.12in}

\begin{extensionbox}
\small
\textbf{Optional Extension --- Compare Two Displays}\\[4pt]
Find a real dataset online with at least 20 quantitative values (e.g., from
\texttt{data.census.gov} or a sports statistics site). Determine the best
display (dotplot, stemplot, or histogram) for your dataset and explain your
reasoning in 3--5 sentences. What does your chosen display reveal that the
other two would not?\\[8pt]
\writeline\\[1pt]\writeline\\[1pt]\writeline\\[1pt]\writeline\\[1pt]\writeline
\vspace{4pt}
\textcolor{slate}{\scriptsize Source / URL:\;\blank{11cm}}
\end{extensionbox}

\end{document}
